%$%%%%%%%%%%%%%%%%%%          First section
\subsection{BB84}
\label{sec:bb84}
\begin{figure}[ht]
    \centering
    \includegraphics[width=0.6\linewidth]{images.methodology/QuantumAbstractTNR.png}
    \caption{Abstract P\&M QKD Network Model}
    \label{fig:QKDnetworkmodel}
\end{figure}

The BB84 protocol, proposed by Bennett and Brassard in 1984 \cite{BB84}, was the first practical QKD scheme. It is a prepare-and-measure protocol that encodes classical bit values onto single-photon polarisation states drawn from two mutually unbiased bases: the rectilinear basis ($\ket{0},\ket{1}$) and the diagonal basis ($\ket{+},\ket{-}$). The sender, Alice, randomly selects both a bit value and a basis for each qubit, producing one of four possible states:
%
\begin{equation}
  \ket{0},\quad \ket{1},\quad
  \ket{+} = \frac{\ket{0}+\ket{1}}{\sqrt{2}},\quad
  \ket{-} = \frac{\ket{0}-\ket{1}}{\sqrt{2}}.
  \label{eq:bb84-states}
\end{equation}
%


These states are transmitted to Bob over the quantum channel. Bob then independently and randomly chooses a measurement basis for each incoming qubit. Because the two bases are mutually unbiased, a measurement in the wrong basis yields a uniformly random outcome and provides no information about Alice's bit value \cite{QKDbible}.

After transmission is complete, Alice and Bob switch to an authenticated classical channel to perform \textit{sifting}, as seen in Fig. \ref{fig:QKDnetworkmodel}. At this point, Bob announces only his choice of basis for each qubit, and Alice confirms which choices matched her preparation basis. All rounds in which the bases disagree are discarded. In an ideal, noiseless channel, the remaining \textit{sifted key} is perfectly correlated between Alice and Bob, and roughly half the transmitted qubits survive the sifting step \cite{BB84}.

\begin{figure}[h]
    \centering
    \includegraphics[width=0.7\linewidth]{images.methodology/bb84processTNR.png}
    \caption{BB84 Operation Model}
    \label{fig:BB84operation}
\end{figure}

To check for eavesdropping, Alice and Bob sacrifice a random subset of the sifted key and compare bit values over the public channel. Any intercept-resend attack by an eavesdropper, Eve, necessarily disturbs a fraction of the qubits because she cannot simultaneously determine the polarisation in both bases without violating the no-cloning theorem \cite{QKDbible}. A full intercept-resend strategy is expected to produce a Quantum Bit Error Rate (QBER) of approximately 25\%, well above the theoretical security threshold of approximately 11\% \cite{ekertQI}. If the observed QBER exceeds this threshold, the protocol aborts. Otherwise, Alice and Bob proceed through error correction and privacy amplification to yield a final secret key, depicted in Fig. \ref{fig:BB84operation}.


\subsection{B92}
\label{sec:b92}

\begin{figure}[h]
    \centering
    \includegraphics[width=0.6\linewidth]{images.methodology/b92basis.png}
    \caption{B92 Measurements Bases}
    \label{fig:B92bases}
\end{figure}

In 1992, Bennett proposed the B92 protocol as a simplification of BB84 \cite{B92}. The central insight is that unconditional security does not require four encoding states; two non-orthogonal states (seen Fig. \ref{fig:B92bases} are sufficient because no measurement can perfectly distinguish between non-orthogonal quantum states \cite{ekertQI}. In B92, Alice encodes bit value 0 as $\ket{0}$ and bit value 1 as $\ket{+}$:

%
\begin{equation}
  \ket{\varphi} =
  \begin{cases}
    \ket{0}  & \text{if } a = 0, \\[4pt]
    \ket{+}  & \text{if } a = 1.
  \end{cases}
  \label{eq:b92-states}
\end{equation}
%

Bob generates his own random bit $a'$ and measures in the $Z$ basis ($\ket{0},\ket{1}$) if $a' = 0$, or in the $X$ basis ($\ket{+},\ket{-}$) if $a' = 1$. When Bob's measurement basis is incompatible with Alice's state, the outcome is conclusive only some of the time, hence, inconclusive results (a.k.a. erasures) are discarded \cite{QKDbible}. Bob publicly announces which rounds yielded a conclusive result, and Alice retains only the corresponding bits. Neither party reveals their basis choice or bit value during this exchange.

Because B92 relies on only two states rather than four, its sifting efficiency is lower than that of BB84, and its tolerance to channel noise is reduced. The theoretical QBER security threshold for B92 is approximately 6.5\%, compared to approximately 11\% for BB84 \cite{ekertQI}. Furthermore, the two-state encoding makes B92 more susceptible to unambiguous state discrimination (USD) attacks, in which Eve exploits the known non-orthogonal structure to occasionally identify Alice's state with certainty while flagging ambiguous cases as losses \cite{B92}. Despite these limitations, B92 remains an important protocol both for its theoretical elegance and because it demonstrates that the minimal resource needed for secure QKD is a pair of non-orthogonal states.

\subsection{E91}
\label{sec:e91}

In contrast to the P\&M paradigm of BB84 and B92, the E91 protocol proposed by Ekert in 1991 \cite{E91} derives its security from quantum entanglement and Bell's theorem. A source prepares pairs of qubits in the singlet Bell state:

%
\begin{equation}
  \ket{\Psi^{-}} = \frac{\ket{01} - \ket{10}}{\sqrt{2}},
  \label{eq:e91-singlet}
\end{equation}
%

and distributes one qubit from each pair to Alice and the other to Bob. Each party independently selects a random measurement orientation from a predetermined set of angles. When Alice and Bob happen to choose the same measurement basis, their outcomes are perfectly anti-correlated, and these matching rounds are used to construct the raw key \cite{E91}.

The rounds in which Alice and Bob chose different bases serve a separate purpose: they are used to evaluate the Clauser-Horne-Shimony-Holt (CHSH) inequality \cite{QKDbible}. For the singlet state, quantum mechanics predicts a CHSH parameter of $S = 2\sqrt{2}$, which maximally violates the classical bound of $S \leq 2$. Any eavesdropping that disturbs the entanglement between the pairs will reduce the observed value of $S$ towards or below the classical limit, thereby revealing Eve's presence \cite{ekertQI, E91}. This connection between eavesdropper detection and Bell inequality violation is the defining feature of E91 and distinguishes it fundamentally from P\&M protocols.

\begin{table}[h]
\centering
\caption{Qualitative Comparison of BB84, B92, and E91}
\label{tab:protocol-comparison}
\footnotesize
\begin{tabularx}{\linewidth}{@{}l X X X@{}}
\toprule
\textbf{Feature} & \textbf{BB84} & \textbf{B92} & \textbf{E91} \\
\midrule
Paradigm & Prepare \& Measure & Prepare \& Measure & Entanglement-Based \\
Encoding states & 4 (two bases) & 2 non-orthogonal & Bell pairs, 3 angles \\
Alice's bases / angles & $\{Z, X\}$ & $\{\ket{0}, \ket{+}\}$ & $\{0, \pi/4, \pi/2\}$ \\
Bob's bases / angles & $\{Z, X\}$ & $\{Z, X\}$ & $\{\pi/4, \pi/2, 3\pi/4\}$ \\
Sifting efficiency & $\sim$50\% & $\sim$25\% & $\sim$22.2\% \\
Threshold ($f_{EC}\!=\!1.0$) & $\sim$11\% & $\sim$8\% & $\sim$11\% \\
Threshold ($f_{EC}\!=\!1.16$) & $\sim$9.8\% & $\sim$6.5\% & $\sim$9.8\% \\
Eve detection & QBER & QBER & CHSH $+$ QBER \\
Full intercept-resend QBER & $\sim$25\% & $\sim$33--38\% & degraded $|S|$ \\
Known vulnerabilities & PNS (WCP sources) & USD attack & Entanglement loss \\
Hardware complexity & Moderate & Lowest & Highest \\
\bottomrule
\end{tabularx}
\end{table}

This entanglement-based approach offers a conceptual advantage as the key bits are not determined until the moment of measurement, so there is no point during transmission at which an encoded classical value exists for Eve to intercept. However, E91 places greater demands on hardware because it requires a source of high-fidelity entangled pairs and the ability to distribute them over the quantum channel while preserving entanglement \cite{ekertQI}. In a simulation context, this translates to the additional complexity of modelling entangled state preparation, Bell-state measurement, and CHSH parameter estimation alongside the standard post-processing pipeline shared with BB84 and B92.



%%%%%%%%%%%%%%%%%%          second section


% >>> ADDED SECTION: Protocol Implementations <<<
\section{Protocol Implementations}
\label{sec:protocol-implementations}

This section describes the implementation of each QKD protocol within the simulation platform. All three protocols share the same post-processing pipeline and noise application mechanism; the differences lie in state preparation, measurement, and sifting logic.

\subsection{BB84 Protocol}

The BB84 implementation follows the four-state, two-basis scheme described in \cref{sec:bb84}. For each of $n$ qubits, Alice randomly selects a bit value $a_i \in \{0,1\}$ and a basis $\hat{a}_i \in \{Z, X\}$ using NumPy's pseudorandom number generator, then prepares the corresponding quantum state on a single-qubit Qiskit circuit. If the selected bit is 1, a Pauli-X gate is applied to flip the qubit from $\ket{0}$ to $\ket{1}$. If the diagonal basis is selected, a Hadamard gate is subsequently applied to rotate into the $X$ basis. An identity gate (\texttt{qc.id(0)}) is appended after state preparation to serve as the channel marker, to which noise models are attached during execution. Bob independently selects a random measurement basis $\hat{b}_i$; if he chooses the diagonal basis, a Hadamard gate is applied before the computational-basis measurement. All $n$ circuits are executed in a single batch call to the Qiskit Aer backend with \texttt{shots=1} and \texttt{memory=True} to retrieve individual outcomes.

The sifting procedure retains only the rounds where $\hat{a}_i = \hat{b}_i$, implemented as a Boolean mask comparison of the two basis arrays. On an ideal channel, this produces a sifting efficiency of approximately 50\% and a QBER of 0\%, serving as the primary validation targets for this protocol. The key rate is computed as the ratio of sifted key length to total transmitted qubits.

% NOTE: Replace firstline/lastline with actual line numbers from bb84.py
\lstinputlisting[
    language=Python,
    caption={BB84 state preparation and channel-marker insertion (\texttt{bb84.py}).},
    label={lst:bb84-prep},
    firstline=1,
    lastline=1
]  {}% TODO: set correct path + line range

The high-level flow of the BB84 key sifting procedure is presented in \cref{alg:bb84}.

\begin{algorithm}[H]
    \caption{BB84 Key Sifting Procedure}\label{alg:bb84}
    \KwIn{Alice's bits $a_i$, bases $\hat{a}_i$; Bob's measurements $b_i$, bases $\hat{b}_i$}
    \KwOut{Sifted key $K$}
    $K \leftarrow \emptyset$\;
    \For{$i \leftarrow 1$ \KwTo $n$}{
        Alice publishes $\hat{a}_i$ over classical channel\;
        Bob publishes $\hat{b}_i$ over classical channel\;
        \If{$\hat{a}_i = \hat{b}_i$}{
            Append $b_i$ to $K$ \Comment{Matching basis -- keep bit}
        } 
    }
    \Return{$K$}
\end{algorithm}


\subsection{B92 Protocol}

The B92 implementation uses only two non-orthogonal encoding states: $\ket{0}$ for bit value 0 and $\ket{+}$ for bit value 1. Alice prepares $\ket{0}$ directly (no gate applied) or applies a Hadamard gate to produce $\ket{+}$, followed by the identity gate channel marker (\texttt{qc.id(0)}).

Bob's measurement logic is the most distinctive aspect of the B92 implementation. Bob generates a random bit $b_i \in \{0, 1\}$ and measures in the $Z$ basis if $b_i = 0$ or the $X$ basis if $b_i = 1$, applying a Hadamard gate before measurement in the latter case. An outcome is \textit{conclusive} only if, given the basis chosen, the observed result is inconsistent with one of Alice's two possible preparations. Measuring $\ket{0}$ in the $Z$ basis can only yield $0$ and is therefore ambiguous, as is measuring $\ket{+}$ in the $X$ basis (which also yields $0$ under the Hadamard-then-$Z$ implementation). The conclusive events are the \emph{complementary} pairings: a $Z$-basis outcome of $1$ rules out a $\ket{0}$ preparation and therefore identifies Alice's bit as $1$; an $X$-basis outcome of $1$ (equivalent to projecting onto $\ket{-}$) rules out $\ket{+}$ and identifies Alice's bit as $0$. Inconclusive rounds (outcome $0$) are discarded.

The sifting step in B92 therefore differs from BB84: Bob announces which rounds were conclusive, not which basis he used. Neither party reveals their bit value or basis choice during this exchange. Bob's sifted key is inferred as $1 - b_i$ (the complement of Bob's measurement basis for the conclusive rounds), while Alice retains her original bit values for those rounds. This sifting mechanism produces a theoretical sifting efficiency of approximately 25\%, half that of BB84. The reduced efficiency, combined with the two-state encoding, results in a lower noise tolerance.

\begin{figure}[ht]
    \centering
    \includegraphics[width=0.6\linewidth]{images.methodology/b92 flowchart.png}
    \caption{B92 algorithm flowchart}
    \label{fig:pipeline}
\end{figure}

% NOTE: Replace firstline/lastline with actual line numbers from b92.py
\lstinputlisting[
    language=Python,
    caption={B92 conclusive-measurement decoding and sifting (\texttt{b92.py}).},
    label={lst:b92-sift},
    firstline=1,
    lastline=1
]{}  % TODO: set correct path + line range


\subsection{E91 Protocol}

The E91 implementation models the entanglement-based protocol described in \cref{sec:e91}. For each round, a two-qubit Qiskit circuit is initialised and a singlet Bell state $\ket{\Psi^{-}} = (\ket{01} - \ket{10})/\sqrt{2}$ is prepared by applying the following gate sequence: a Hadamard gate on qubit 0, a CNOT gate with qubit 0 as control and qubit 1 as target, a Pauli-X gate on qubit 0, and a Pauli-Z gate on qubit 0. This sequence produces the singlet state rather than the more commonly used $\ket{\Phi^+}$, matching the original E91 specification \cite{E91}. Qubit 0 is designated as Alice's and qubit 1 as Bob's.

The implementation supports two channel topologies via a configurable parameter. In the \texttt{bob} topology, an identity gate channel marker is applied only to Bob's qubit (\texttt{qc.id(1)}), representing a scenario where noise acts on the distribution channel to Bob while Alice's qubit remains local. In the \texttt{both} topology, identity gates are applied to both qubits (\texttt{qc.id(0)} and \texttt{qc.id(1)}), modelling noise on both distribution paths. This distinction affects the relationship between noise strength and observed QBER: under depolarising noise with parameter $p$, the \texttt{bob} topology produces $\text{QBER} = p/2$ (matching BB84), while the \texttt{both} topology produces $\text{QBER} = p - p^2/2$.

\begin{figure}[ht]
    \centering
    \includegraphics[width=0.6\linewidth]{images.methodology/e91 bloch sphere.png}
    \caption{E91 XZ plane measurement bases}
    \label{fig:pipeline}
\end{figure}

Alice and Bob each independently select a measurement angle from predetermined sets. Alice's angles are drawn from $\{0, \pi/4, \pi/2\}$ and Bob's from $\{\pi/4, \pi/2, 3\pi/4\}$, each chosen uniformly at random per round. Measurement in an arbitrary angle $\theta$ is implemented by applying a rotation gate $R_y(-\theta)$ before a $Z$-basis measurement. A predefined $3 \times 3$ sift table, indexed by Alice's and Bob's angle choices, classifies each round as contributing to the key, to the CHSH parameter estimation, or as discarded. Of the nine possible angle-pair combinations, two contribute to the key (matching angles that yield anti-correlated outcomes), four contribute to CHSH estimation, and three are discarded, producing a theoretical sifting efficiency of $2/9 \approx 22.2\%$.

The CHSH parameter is computed as
%
\begin{equation}
\label{eq:chsh-computation}
S = +E(a_1, b_1) - E(a_1, b_3) + E(a_3, b_1) + E(a_3, b_3),
\end{equation}
%
where $E(a_i, b_j)$ is the correlation coefficient for the angle pair $(a_i, b_j)$, calculated as the expectation value of the product of Alice's and Bob's outcomes (mapped to $\pm 1$ via $1 - 2 \times \text{outcome}$) across all rounds using that pair. The implementation is validated by confirming that $|S| \rightarrow 2\sqrt{2} \approx 2.828$ in the absence of noise or eavesdropping, consistent with the maximal quantum violation of the CHSH inequality. Under depolarising noise, the theoretical CHSH value is $2\sqrt{2}(1-p)$ for the \texttt{bob} topology and $2\sqrt{2}(1-p)^2$ for the \texttt{both} topology.

% NOTE: Replace firstline/lastline with actual line numbers from e91.py
\lstinputlisting[
    language=Python,
    caption={E91 Bell-state preparation and CHSH computation (\texttt{e91.py}).},
    label={lst:e91-chsh},
    firstline=1,
    lastline=1
]{}  % TODO: set correct path + line range
% >>> END ADDED SECTION <<<
